\documentclass[10pt]{article}
\usepackage[T1]{fontenc}
\usepackage[default]{sourcesanspro}
\usepackage[letterpaper,top=0.85in,bottom=0.8in,left=0.7in,right=0.7in,
  headheight=14pt,headsep=0.25in,footskip=0.4in]{geometry}
\usepackage{fancyhdr}
\usepackage{xcolor}
\usepackage{colortbl}
\usepackage{array}
\usepackage{booktabs}
\usepackage{tabularx}
\usepackage{enumitem}
\usepackage{titlesec}
\usepackage{needspace}
\usepackage{microtype}
\usepackage[italic]{mathastext}

\definecolor{rule}{gray}{0.55}
\definecolor{head}{gray}{0.90}
\definecolor{box}{gray}{0.95}

\setlength{\parindent}{0pt}
\setlength{\parskip}{4pt}
\linespread{1.04}

\pagestyle{fancy}
\fancyhf{}
\fancyhead[L]{Week 7 / Take-away games / Adult guide}
\fancyfoot[L]{Bellingham Math Circle / Week 7 / F07-FAC-v4}
\fancyfoot[R]{\thepage}
\renewcommand{\headrulewidth}{0pt}
\renewcommand{\footrulewidth}{0pt}

\titleformat{\section}{\large\bfseries}{\thesection}{0.6em}{}
\titlespacing*{\section}{0pt}{12pt}{4pt}
\titleformat{\subsection}{\normalsize\bfseries}{\thesubsection}{0.5em}{}
\titlespacing*{\subsection}{0pt}{9pt}{3pt}

\setlist{nosep,leftmargin=1.2em,topsep=2pt}
\setlist[itemize]{label=\textbullet}

% table helpers
\newcolumntype{L}[1]{>{\raggedright\arraybackslash}p{#1}}
\newcolumntype{Y}{>{\raggedright\arraybackslash}X}
\renewcommand{\arraystretch}{1.12}
\setlength{\tabcolsep}{4pt}
\newcommand{\hd}{\rowcolor{head}}
% problem tables: one small table per problem, so page breaks fall between
% problems
\newlength{\cA}\newlength{\cB}
\newenvironment{probtable}[2]{\par\setlength{\cA}{#1}\setlength{\cB}{#2}%
  \setlength{\parskip}{0pt}\noindent
  \begin{tabularx}{\linewidth}{@{}L{\cA}L{\cB}Y@{}}\toprule
  \hd Problem and what they do & Answer & If stuck, in this order \\ \midrule
  \end{tabularx}\par\nopagebreak}{\par}
\newcommand{\prow}[3]{\noindent
  \begin{tabularx}{\linewidth}{@{}L{\cA}L{\cB}Y@{}}#1 & #2 & #3 \\ \midrule
  \end{tabularx}\par}
% numbered hint lines inside a cell
\newcommand{\hi}[1]{\par\noindent\hangindent=0.9em\hangafter=1\makebox[0.9em][l]{#1}}
% something to say
\newcommand{\say}[1]{\par\textit{Say:} \textit{``#1''}}
\newcommand{\core}{\textsc{core}}
\newcommand{\further}{\textsc{further}}
\newcommand{\tagc}{\par{\footnotesize\textbf{core}}}
\newcommand{\tagf}{\par{\footnotesize\textbf{further}}}
\newcommand{\ra}{$\to$}
\newcommand{\shaded}[1]{\par\noindent\colorbox{box}{\parbox{\dimexpr\linewidth-2\fboxsep}{#1}}\par}

\begin{document}

\section{The idea of the week}

Two players take turns taking counters from a pile. Each turn takes an allowed
number (for example 1 or 2), and whoever takes the last counter wins. The same
game is played by moving one token down a number track toward 0 by the number
taken; whoever puts the token on 0 wins.

Every pile size is either one you want to hand to your opponent or one you
want to be handed. Working up from 0 sorts them all: a pile is one you want to
hand over when every move from it leads to a pile you do not want to hand
over. With 1 or 2 those piles are 0, 3, 6, 9, \dots{} The pattern always
repeats, and it changes when the rule changes. Taking as many as you can is
often wrong. A winning strategy must answer every reply, not only the one you
expect.

\textbf{What a good session looks like}

\begin{itemize}
\item \textbf{K--1.} Every child plays many short games with real counters and
with the token, loses a few to the grown-up, then says the grown-up's secret
(``land on 9, 6, 3'') and uses it to win against a partner. Quick children see
that the secret changes when the rule changes.
\item \textbf{2--3.} Pairs play the 1, 3 or 4 game until they choose first or
second correctly, colour the squares they want to leave (0, 2, 7, 9, 14, 16),
and test the claims about Lena and Omar against every reply, not one.
\item \textbf{4--5.} The children play each rule before colouring it, find
the repeating pattern for several rules, and explain why the 1, 3 or 4 pattern
repeats every 7. The quickest reach the last-counter-loses and two-pile games.
\end{itemize}

\section{Materials and preparation}

\noindent\begin{tabularx}{\linewidth}{@{}L{1.35in}YYYL{0.75in}@{}}
\toprule
\hd Item & K--1: 4 children, 2 pairs & 2--3: 4 children, 2 pairs & 4--5: 3 children & Total \\
\midrule
Counters, one colour &
12 per pair (largest pile 10) + 6 spare = 30 &
12 per pair (largest 5 + 6 = 11) + 6 spare = 30 &
20 for the pair (largest 16) + 12 to build 20 and 11 + 4 spare = 36 &
96, plus 7 for the launch \\
Tokens for the track &
1 per pair + 1 spare = 3 & 1 per pair + 1 spare = 3 & 1 + 1 spare = 2 &
8, plus 1 for the launch \\
Coloured pencils, light colours &
2 per child = 8 & 2 per child = 8 & 2 per child = 6 & 22 \\
Pencils & 4 + 1 adult = 5 & 4 + 1 = 5 & 3 + 1 = 4 & 14 \\
Small whiteboards and markers & none & 1 per pair + 1 adult = 3 & 1 per child + 1 adult = 4 & 7 \\
Blank paper & none & 4 sheets & 6 sheets & 10 \\
\bottomrule
\end{tabularx}

\textbf{Tokens.} The track squares are about 0.68 inch (17\,mm). Use a small
token: a snap cube or a counter of a different colour from the pile counters.
A pattern block is too big.

\textbf{Printing.} US Letter, single-sided, at 100\% (``Actual size'', not
``Fit to page''); the track squares must come out about 0.68 inch. Each child
gets their own copy and one page at a time: hand out the next page when a
child finishes one. Spare track pages are for children who want to play again.

\noindent\begin{tabularx}{\linewidth}{@{}L{1.2in}L{1.3in}Y@{}}
\toprule
\hd Table & Pages & Copies \\
\midrule
K--1 & 1--6 (Problems 1--9) & 6 of each page (4 + 2 spare) = 36 sheets, plus 4 more of page 2 (six 0--10 tracks) \\
2--3 & 1--5 (Problems 1--9) & 6 of each page = 30 sheets, plus 2 more of page 1 (two 0--20 tracks) \\
4--5 & 1--8 (Problems 1--14) & 5 of each page (3 + 2 spare) = 40 sheets \\
Adults & this guide; the packet for their table & 3 guides; 1 packet per table \\
Launch & K--1 page 2 & 1 extra copy (its 0--10 track) \\
\bottomrule
\end{tabularx}

\textbf{Set up.} Nothing to cut or tape. Before the children arrive, put at
each table a cup with that table's counters (counted), the tokens, the
pencils and coloured pencils, and the pages face down in order with page 1 on
top. Put 7 counters, a token and the extra K--1 page 2 at the table where the
launch happens.

\needspace{5\baselineskip}
\textbf{Test with the real materials (five minutes).}
\begin{itemize}
\item Put your token on a printed track. It should sit on one square without
hiding the numbers next to it. (A 2\,cm snap cube may hide the number it sits
on; that is fine.)
\item The real counters are bigger than the drawn circles (0.4 inch). Children
build each pile next to the picture, not on it. Check that a pile of 10 fits
in front of a pair.
\item Colour one square with each coloured pencil; the number must still be
readable. Leave out dark colours.
\item Play your table's first rule once against yourself using the tables in
Section 4, so the winning moves are familiar.
\end{itemize}

\textbf{Tears.} Competitive games can bring tears. Keep rounds short (a pile
of 10 takes under a minute), swap who goes first every game, and start the next
game at once instead of dwelling on a loss. At 4--5 the roles rotate every
game. A child who is upset chooses the pile and who starts in the next game.

\section{The hour}

\noindent\begin{tabularx}{\linewidth}{@{}L{0.95in}Y@{}}
\toprule
\hd Time & What happens \\
\midrule
0:00--0:05 & Running around. \\
0:05--0:08 & Whole-group launch (organizer), script below. \\
0:08--0:48 & Work at the three tables in pairs. \\
0:48--0:55 & Sharing: one question, below. \\
0:55--1:00 & Counters back in the cups (count them); pages collected. \\
\bottomrule
\end{tabularx}

\textbf{Launch, two to three minutes.} Everyone stands around one table with
7 counters in a row and the token on square 7 of a 0--10 track.
\begin{enumerate}[leftmargin=1.4em]
\item Point to the counters. \textit{``Here are 7 counters. Two players take
turns. On your turn you take 1 counter or 2 counters. Whoever takes the last
counter wins.''}
\item Point to the track. \textit{``The token shows how many are left. When
you take counters, move the token that many squares toward 0.''}
\item \textit{``Who will play me? I'll go first.''} Take 1 counter and move the
token from 7 to 6. \textit{``I took 1. Six left, and the token is on 6.''}
\item The child takes 1 or 2 counters and moves the token the same number of
squares. Ask the group: \textit{``Is that allowed?''} If a child tries to take
3: \textit{``Only 1 or 2. Put one back.''}
\item You take whatever makes 3 with the child's take, so the token lands on
3. The child moves; you take the rest. \textit{``I took the last counter, so I
win.''}
\item \textit{``The grown-ups know a secret way to win. Your job today is to
beat the grown-up and find the secret. Some tables change the rule, so listen
to your page.''} Send the children to their tables.
\end{enumerate}

\textbf{Pairs.}
\begin{itemize}
\item \textbf{K--1}: two pairs, the two kindergartners together and the two
first graders together, so games are even. Each pair shares a cup of counters
and a token; each child has a page. When the page says ``play the
grown-up'', play one pair as a team (the two children take turns making the
team's move) while the other pair keeps playing the problem before; then swap.
\item \textbf{2--3}: two pairs. For Problem 3 the adult plays one pair while
the other works on Problem 2 or 4.
\item \textbf{4--5}: two players and a referee, switching after every game
(A and B with C refereeing, then B and C, then C and A). The referee checks
that each move is legal, says whose turn it is and keeps count. Instead of
refereeing, the third child may play the adult: choose the pile and who goes
first, and try to win. The organizer plays whoever is waiting.
\end{itemize}

\textbf{When attention runs out.} A standing-up counting game; each is a
take-away game in disguise, and the adult plays to win.
\begin{itemize}
\item \textbf{K--1, Race to 10.} Count aloud together from 1. On your turn
say one or two numbers. Whoever says 10 wins. The grown-up goes first and says
1, then makes sure to say 4, 7 and 10.
\item \textbf{2--3, Race to 21.} Say one, two or three numbers. The adult goes
first and says 1, then 5, 9, 13, 17 and 21.
\item \textbf{4--5, Race to 100.} Add any whole number from 1 to 10 to the
running total. Whoever reaches exactly 100 wins. The organizer goes first and
says 1, then 12, 23, 34, 45, 56, 67, 78, 89 and 100.
\end{itemize}

\textbf{Sharing question} (organizer, to everyone): \textit{``Tell us a pile
where you would rather go second, and what the rule was.''} Write the answers on
the board in one column per rule.

\section{Table by table}

\subsection*{How the adult wins, every rule on the pages}

Play to win, and let the children hunt for the secret. If you are handed one
of the squares you want to leave, there is no good move: take 1 and wait for
a mistake. ``Remainder'' means the remainder on dividing by the number in the
second column.

\noindent\begin{tabularx}{\linewidth}{@{}L{1.05in}L{1.45in}L{2.15in}Y@{}}
\toprule
\hd Rule & Where & Squares to leave & From any other pile, take \\
\midrule
1 or 2 & Launch; K--1 P1--5; 4--5 P1, P2, P6 & 0, 3, 6, 9, 12, \dots{} (multiples of 3) & what makes 3 with their take: they 1, you 2; they 2, you 1 \\
1, 2 or 3 & K--1 P6; 4--5 P2, P6 & 0, 4, 8, 12, 16, 20 (multiples of 4) & what makes 4 with their take \\
1 or 4 & K--1 P7; 2--3 P8, P9; 4--5 P7 & 0, 2, 5, 7, 10, 12, 15, 17, \dots{} (remainder 0 or 2, dividing by 5) & remainder 1: take 1; 3: take 1; 4: take 4 \\
1, 3 or 4 & 2--3 P1--7; 4--5 P3--6, P8 & 0, 2, 7, 9, 14, 16, 21, 23, 28, 30, \dots{} (remainder 0 or 2, dividing by 7) & remainder 1: take 1; 3: take 1 or 3; 4: take 4; 5: take 3; 6: take 4 \\
2 or 3 & 4--5 P7 & 0, 1, 5, 6, 10, 11, \dots{} (remainder 0 or 1, dividing by 5) & remainder 2: take 2; 3: take 2 or 3; 4: take 3 \\
1, 3 or 5 & 2--3 P9; 4--5 P7 & the even squares & take 1 \\
Last counter loses, 1 or 2 & K--1 P9; 4--5 P10, P11 & 1, 4, 7, 10, 13, 16, 19 & remainder 2: take 1; remainder 0: take 2 \\
Last counter loses, 1, 2 or 3 & 4--5 P11 & 1, 5, 9, 13, 17 & remainder 2: take 1; 3: take 2; 0: take 3 \\
Two piles, take from one & K--1 P8; 2--3 P7; 4--5 P12, P13 & equal piles; with 1 or 2, any two piles with the same remainder dividing by 3 & from equal piles: copy their move in the other pile \\
2, 5 or 6 & 4--5 P14 (colouring only) & 0, 1, 4, 8, 11, 12, 15, 19, 22, 23, 26, 30 & \\
\bottomrule
\end{tabularx}

\needspace{6\baselineskip}
\subsection{K--1 table (parent volunteer)}

\shaded{\textbf{The grown-up's secret with 1 or 2:} land on 9, 6, 3, then 0.
Whatever the child takes, you take enough to make 3 together: they take 1, you
take 2; they take 2, you take 1. If the token is already on 3, 6 or 9 when it
is your turn, there is no good move: take 1 and hope. Read each problem aloud
once and point to the picture. Do not tell the secret; say \textit{``Watch
where I put the token.''}}

\textbf{Core:} Problems 1--5. \textbf{Further:} 6--9. \textbf{If a child is
struggling:} in Problem 1 do piles 2 to 5 only, then Problems 2 and 3; skip 4
and 5.

\begin{probtable}{1.85in}{2.2in}
\prow{\textbf{P1} \textit{(core)} Make piles of 2 to 8 counters, play each many times with the
partner, circle 1st or 2nd.
\say{Play it again, and swap who goes first.}}
{Go \textbf{2nd} with 3 and 6. Go \textbf{1st} with the others: from 2 take 2;
from 4 take 1; from 5 take 2; from 7 take 1; from 8 take 2 (always leave 3 or
6, or take the last).}
{\hi{1}Make 3 counters. ``You go first. What can you take? What does your
partner do then?''
\hi{2}``Would you rather hand your partner 2 counters or 3?''
\hi{3}``From 5, can you leave your partner 3?''}
\prow{\textbf{P2} \textit{(core)} The pair plays the grown-up from 10 as a team; the grown-up goes
first. They colour each square where the grown-up puts the token. Three
games.
\say{I'm going first, and I'm going to try to win. Colour my square every
time I move. Can you find my secret?}}
{You take 1 (10 \ra{} 9), then make 3 with their take each time. You land on
9, 6, 3, 0 in every game, so the same four squares get coloured on all three
tracks.}
{\hi{1}``Look at your three tracks. Which squares are coloured on all of them?''
\hi{2}``Where did I put the token first, every time?''
\hi{3}``If the token is on 3 and it is your turn, can you win?''}
\prow{\textbf{P3} \textit{(core)} The child goes first from 10 against the partner and colours
their own squares.}
{The first player wins by landing on 9, 6, 3, 0 (take 1 first). A child who
colours exactly 9, 6, 3, 0 and wins has found the secret.}
{\hi{1}``Where did I go first in Problem 2?''
\hi{2}``Put the token on 9. What can your partner do now?''
\hi{3}``Your partner moved to 7. Can you get to 6?''}
\prow{\textbf{P4} \textit{(core)} Piles of 5, 7, 8, 10. The child goes first and circles the
counters they take first so they can win.}
{5: take 2 (leave 3). 7: take 1 (leave 6). 8: take 2 (leave 6). 10: take 1
(leave 9). Each is the only winning first move.}
{\hi{1}``Which piles did you want to give your partner in Problem 1?''
\hi{2}Make the pile of 7. ``How many do you take to leave 6?''
\hi{3}Play it out with the child's choice and see.}
\prow{\textbf{P5} \textit{(core)} The partner puts the token on a square from 1 to 10; the child
circles it and circles 1st or 2nd. Three games.}
{On 3, 6 or 9: \textbf{2nd}. Anywhere else: \textbf{1st}, and move to the next
of 9, 6, 3, 0 below (from 1, 4, 7, 10 take 1; from 2, 5, 8 take 2).}
{\hi{1}``Is the token on one of my secret squares?''
\hi{2}``If it is, do you want to move first?''
\hi{3}``If not, which secret square can you reach?''}
\prow{\textbf{P6} \textit{(further)} Now 1, 2 or 3 squares from 10. First the grown-up goes first
(two games), then the child goes first against the partner (three games).
They colour the first player's squares.}
{The first player lands on 8, 4, 0: take 2 first, then make 4 with their take
(they 1, you 3; they 2, you 2; they 3, you 1). The old secret fails: from 6
the other player takes 2 and lands on 4.}
{\hi{1}``Is the secret still 9, 6, 3? Leave your partner 6 and see.''
\hi{2}``Which squares did the grown-up colour?''
\hi{3}``How far apart are 8, 4 and 0?''}
\prow{\textbf{P7} \textit{(further)} Now 1 or 4 squares from 8, the same way as P6.}
{The first player takes 1 (8 \ra{} 7). The partner goes to 6 or 3. From 6 go
to 5 (take 1) or 2 (take 4); from 3 go to 2. From 5 or 2 the first player
reaches 0 next. So the coloured squares are 7, then 5 or 2, then 0. Taking 4
first (8 \ra{} 4) loses: the other player takes 4 and lands on 0.}
{\hi{1}``Take 4 first. What happens?''
\hi{2}``Where did the grown-up go first?''
\hi{3}``If the token is on 2 and it's your partner's turn, what can they do?''}
\prow{\textbf{P8} \textit{(further)} Two piles; take 1 or 2 from one pile; whoever takes the last
counter on the table wins. Circle 1st or 2nd for 2 and 2, 1 and 2, 3 and 3,
2 and 4, 4 and 4.}
{Equal piles (2 and 2, 3 and 3, 4 and 4): \textbf{2nd}, then copy: whatever
your partner takes from one pile, take the same from the other. 1 and 2:
\textbf{1st}, take 1 from the 2. 2 and 4: \textbf{1st}, take 2 from the 4.
Why copying wins: ``After my turn the piles are the same. Whatever you do, I
can do the same, so I always have a move, and I take the last one.''}
{\hi{1}Make 2 and 2. ``You go first. What happens?''
\hi{2}``What if you do exactly what your partner did, in the other pile?''
\hi{3}``With 1 and 2, can you make the piles the same?''}
\prow{\textbf{P9} \textit{(further)} Now whoever takes the last counter loses. Piles 1 to 6, circle
1st or 2nd.}
{1: \textbf{2nd}. 2: 1st, take 1. 3: 1st, take 2. 4: \textbf{2nd}. 5: 1st,
take 1. 6: 1st, take 2. Leave your partner 1 or 4.}
{\hi{1}``With 1 counter, who has to take it?''
\hi{2}``Can you leave your partner just 1?''
\hi{3}``Which pile makes your partner end up with 1?''}
\end{probtable}

\needspace{6\baselineskip}
\subsection{2--3 table (mathematician)}

\shaded{Rule for P1--P7: take 1, 3 or 4. Squares to leave: 0, 2, 7, 9, 14,
16, 21, 23, 28, 30, \dots{} (remainder 0 or 2 on dividing by 7). Replies are in
the table at the start of this section.}

\textbf{Core:} Problems 1--5. \textbf{Further:} 6--9. \textbf{If a child is
struggling:} Problem 1 with piles 2, 5 and 7 only, then play the adult
(Problem 3); skip Problems 2 and 4.

\begin{probtable}{1.7in}{2.45in}
\prow{\textbf{P1} \textit{(core)} Piles 2, 5, 6, 7, 8, moves 1, 3 or 4; circle first or second
and the first take.}
{2: \textbf{second}. 5: first, take 3. 6: first, take 4. 7:
\textbf{second}. 8: first, take 1. Each winning first move is the only one.}
{\hi{1}``With 2 counters, what can the first player take?''
\hi{2}``Which small piles do you want to hand your partner?''
\hi{3}``From 7, try each move (to 6, 4, 3). Can your partner always leave you
2 or 0?''}
\prow{\textbf{P2} \textit{(core)} Colour every square of 0--20 they would like to leave. Second
track for a fresh try.}
{0, 2, 7, 9, 14, 16.}
{\hi{1}``Is 1 a square you want to leave? Is 2?''
\hi{2}``Go up one square at a time. Can your opponent move from this square
to a coloured one? Then don't colour it.''
\hi{3}``Check 7: your opponent can go to 6, 4 or 3. Can you reach a coloured
square from each?''}
\prow{\textbf{P3} \textit{(core)} The pair plays the adult. The adult puts the token on 10--20;
the pair chooses who goes first, circles the start and colours their
squares.}
{Go second from 14 or 16. Otherwise go first and move to: from 10, 9 or 7;
11 \ra{} 7; 12 \ra{} 9; 13 \ra{} 9; 15 \ra{} 14; 17 \ra{} 16 or 14; 18 \ra{}
14; 19 \ra{} 16; 20 \ra{} 16. Use 14 and 16 early, then starts like 13 and
18 where only one move wins. If they choose wrongly, win with the reply
table.}
{\hi{1}``Is the start one of your coloured squares?''
\hi{2}``If the start is 14, do you want to move first?''
\hi{3}``From 18, which coloured square can you reach?''}
\prow{\textbf{P4} \textit{(core)} Start 9. Lena says that if her opponent goes first she wins
whatever the opponent does. Right? Explain.}
{\textbf{Yes.} The opponent can go to 8, 6 or 5; Lena answers 8 \ra{} 7, 6
\ra{} 2, 5 \ra{} 2. From 7 the opponent goes to 6, 4 or 3; Lena answers 6
\ra{} 2, 4 \ra{} 0, 3 \ra{} 0. From 2 the opponent must take 1; Lena takes
the last. A child's explanation: ``Whatever you do from 9, I can get to 7 or
2. From 7, whatever you do, I get to 2 or 0. From 2 you can only take 1.''}
{\hi{1}``Be Lena. I'll go first from 9.'' Then: ``Now I'll try a different
first move.''
\hi{2}``How many different first moves does the opponent have?''
\hi{3}``For each one, where can Lena go?''}
\prow{\textbf{P5} \textit{(core)} Omar always takes as many as he is allowed and goes first from
13. Does he win whatever the opponent does? Then find a way that does.}
{\textbf{No.} His first move (take 4, to 9) is right, but later ones go wrong.
Opponent to 8: Omar takes 4 (to 4), opponent takes 4 and wins. Opponent to 5:
Omar takes 4 (to 1), opponent takes the last. (Opponent to 6: Omar wins.)
Winning way: take 4 to 9, then play Lena's answers from P4.}
{\hi{1}``Play Omar's way from 13. I'll be the opponent.''
\hi{2}``After Omar goes to 9, try each move for the opponent.''
\hi{3}``From 8, where should Omar go?'' (7)}
\prow{\textbf{P6} \textit{(further)} Piles of 50 and 100: first or second, and how many to take.
Explain.}
{50: \textbf{first, take 1} (to 49). 100: \textbf{second}. The coloured
squares come in pairs 7 apart: 0, 2, 7, 9, 14, 16, \dots{} Each multiple of 7
and the square 2 above it. 49 = 7 $\times$ 7, so 49 and 51 are coloured and 50
is not; 98 = 7 $\times$ 14, so 98 and 100 are coloured. Why it keeps going: a
square only looks at the squares 1, 3 and 4 below it, and 7 to 10 look like 0
to 3, so the pattern copies itself every 7.}
{\hi{1}``Colour the 0--30 track. How far apart are the coloured squares?''
\hi{2}On a whiteboard write 0 to 34 in rows of 7 and circle the coloured ones
(two columns appear).
\hi{3}``Count by 7s up to 50.''}
\prow{\textbf{P7} \textit{(further)} Two piles, 1, 3 or 4 from one pile: 5 and 5, then 5 and 6.
First or second? Explain why it wins whatever the opponent does.}
{5 and 5: \textbf{second}, copy the opponent in the other pile. 5 and 6:
\textbf{first}, take 1 from the 6, then copy. Why: ``After my turn the piles
are equal. Whatever you take, I can take the same from the other pile, so I
always have a move and I take the last counter.'' (Taking 1 from the 5, to 4
and 6, also wins; see Section 5.)}
{\hi{1}``Play 5 and 5. Whatever I do in one pile, what could you do in the
other?''
\hi{2}``If the piles are equal after your turn, can your opponent take the
last counter?''
\hi{3}``With 5 and 6, can you make them equal?''}
\prow{\textbf{P8} \textit{(further)} Moves 1 or 4; piles 5, 6, 8, 9, 10: first or second and the
first take.}
{5: \textbf{second}. 6: first, take 1 or 4. 8: first, take 1. 9: first, take
4. 10: \textbf{second}.}
{\hi{1}``Which small piles do you want to hand over? Try 1, 2 and 3.''
\hi{2}``From 5, where can your opponent go?'' (4 or 1; you take all.)
\hi{3}``Colour a 0--10 track for this rule first.''}
\prow{\textbf{P9} \textit{(further)} Colour the squares to leave for 1 or 4, and for 1, 3 or 5.}
{1 or 4: 0, 2, 5, 7, 10, 12, 15, 17, 20 (repeats every 5). 1, 3 or 5: every
even square, 0 to 20. Why: every move is odd, so from an even square you
always land on an odd one, and from an odd square you take 1.}
{\hi{1}``Which of 0, 1, 2 do you colour?''
\hi{2}``Compare with Problem 8.''
\hi{3}For 1, 3, 5: ``What happens to odd and even when you take an odd
number?''}
\end{probtable}

\needspace{6\baselineskip}
\subsection{4--5 table (organizer)}

\textbf{Core:} Problems 1--6. \textbf{Further:} 7--14 (7--9 change the rule and
explain the repetition; 10--11 last counter loses; 12--13 two piles; 14
pigeonhole). \textbf{If a child is struggling:} play more games of Problems 1
and 3 with the referee, and leave the colouring (2 and 4) and Kai (5) until
the games are settled.

\begin{probtable}{1.6in}{2.75in}
\prow{\textbf{P1} \textit{(core)} Moves 1 or 2; piles 7, 9, 11, 12, 16.}
{7: first, take 1. 9: \textbf{second}. 11: first, take 2. 12:
\textbf{second}. 16: first, take 1.}
{\hi{1}``Which pile would you hand over: 1, 2 or 3?''
\hi{2}``Play 6 a few times. Can the first player win against a careful
second player?''
\hi{3}``What do 3, 6, 9 and 12 have in common?''}
\prow{\textbf{P2} \textit{(core)} Colour squares to leave, 0--20, for 1 or 2 and for 1, 2 or 3.}
{1 or 2: 0, 3, 6, 9, 12, 15, 18. 1, 2 or 3: 0, 4, 8, 12, 16, 20. Strategy:
answer a take of $k$ with $3-k$ (or $4-k$).}
{\hi{1}``Start at 0. Is 1 a square you want to leave? 2? 3?''
\hi{2}``Colour a square when every move from it lands on an uncoloured one.''
\hi{3}``If your opponent takes 1, what do you take so the round removes 4?''}
\prow{\textbf{P3} \textit{(core)} Moves 1, 3 or 4; piles 5, 8, 9, 11, 14.}
{5: first, take 3. 8: first, take 1. 9: \textbf{second}. 11: first, take 4. 14:
\textbf{second}.}
{\hi{1}``Is a pile of 2 good to hand over?''
\hi{2}``Work up from 0 to 14 on the Problem 4 track first.''
\hi{3}``From 11, which of 0, 2, 7, 9 can you reach?''}
\prow{\textbf{P4} \textit{(core)} Squares to leave for 1, 3 or 4, 0--30.}
{0, 2, 7, 9, 14, 16, 21, 23, 28, 30 (remainder 0 or 2 on dividing by 7).}
{\hi{1}``Which of the piles in Problem 3 did you want to hand over?''
\hi{2}``For each square, look 1, 3 and 4 squares down.''
\hi{3}Write 0--34 in rows of 7 on a whiteboard.}
\prow{\textbf{P5} \textit{(core)} Kai goes first and takes as many as allowed on his first turn.
Colour every pile 1--30 that the first player can win but Kai's first move
spoils.}
{5, 8, 10, 12, 15, 17, 19, 22, 24, 26, 29: remainder 1, 3 or 5 on dividing by
7, from 5 up. From $n \ge 4$ Kai leaves $n-4$, a pile to leave only when $n$
has remainder 4 or 6. If a child reads it as greedy for the whole game, add
13, 18, 20, 25, 27 (16 piles; e.g.\ 13 \ra{} 9 \ra{} 8 \ra{} 4 \ra{} 0).
Greedy all game wins against every reply only from 1, 3, 4, 6, 11.}
{\hi{1}``From 12, where does Kai's first move go? Is that square coloured?''
\hi{2}``Kai always lands 4 below the start. Which starts land on a coloured
square?''
\hi{3}``Slide the coloured squares up by 4.''}
\prow{\textbf{P6} \textit{(core)} Pile of 100 with 1 or 2, with 1, 2 or 3, with 1, 3 or 4: first
or second, first take, and why it wins against every reply.}
{1 or 2: \textbf{first, take 1} (to 99), then answer $k$ with $3-k$; the pile
goes 99, 96, \dots, 3, 0. 1, 2 or 3: \textbf{second}; answer $k$ with $4-k$;
100 = 25 $\times$ 4. 1, 3 or 4: \textbf{second}, since 100 = 98 + 2; keep the
pile at remainder 0 or 2 (dividing by 7) with the reply table at the start of
this section.}
{\hi{1}``Play the same rule with 10 counters first.''
\hi{2}``If your opponent takes 1, what do you take so the two of you took 3?''
\hi{3}For 1, 3, 4: ``Which coloured squares are near 100?'' (91, 93, 98, 100)}
\prow{\textbf{P7} \textit{(further)} Squares to leave, 0--30, for 1 or 4; 2 or 3; 1, 3 or 5.}
{1 or 4: 0, 2, 5, 7, 10, 12, 15, 17, 20, 22, 25, 27, 30. 2 or 3: 0, 1, 5, 6, 10,
11, 15, 16, 20, 21, 25, 26, 30 (1 is coloured because no move is possible from
it). 1, 3 or 5: the even squares.}
{\hi{1}``With 2 or 3, who wins from a pile of 1?''
\hi{2}``Work up from 0.''
\hi{3}``How long before each pattern repeats?''}
\prow{\textbf{P8} \textit{(further)} With 1, 3 or 4, does the pattern repeat forever? Explain.}
{\textbf{Yes, every 7.} A square's colour depends only on the squares 1, 3 and
4 below it. Squares 7, 8, 9, 10 are coloured like 0, 1, 2, 3 (yes, no, yes,
no), so 11 looks at 10, 8, 7, which look like 3, 1, 0; 11 is coloured like 4,
and so on forever. Or by remainders: from 0 or 2 every move lands on 1, 3, 4,
5 or 6, and from those take 1, 1, 4, 3, 4 to get back.}
{\hi{1}``Compare squares 0--6 with 7--13.''
\hi{2}``To decide square 11, which squares do you look at?''
\hi{3}``Those look like 3, 1, 0. Which square used 3, 1, 0?''}
\prow{\textbf{P9} \textit{(further)} Choose moves so the squares to leave are 0, 5, 10, \dots; then
as many choices as possible for 0, 3, 6, \dots; explain.}
{0, 5, 10: moves 1, 2, 3, 4 (exactly the lists containing 1, 2, 3, 4 and no
multiple of 5). 0, 3, 6: exactly the lists containing 1 and 2 and no multiple
of 3: 1, 2; 1, 2, 4; 1, 2, 5; 1, 2, 4, 5; 1, 2, 7; 1, 2, 4, 7; 1, 2, 8;
infinitely many. Why: piles 1 and 2 must be wins, and from them the only
coloured square in reach is 0, so 1 and 2 must be allowed. A move of 3, 6,
\dots{} would jump from a coloured square to a coloured one. With those two
conditions, from any other square take 1 or 2 down to the coloured square.}
{\hi{1}``Try moves 1, 2, 3, 4. Which squares are coloured?''
\hi{2}``For 0, 3, 6: does 1, 2 work? 1, 2, 4? 1, 2, 3?''
\hi{3}``Which moves would let you jump from one coloured square to
another?''}
\prow{\textbf{P10} \textit{(further)} Last counter loses, 1 or 2; piles 1, 3, 5, 7, 8.}
{1: \textbf{second}. 3: first, take 2. 5: first, take 1. 7: \textbf{second}. 8:
first, take 1. Leave 1, 4, 7, \dots}
{\hi{1}``With a pile of 1, who has to take the last counter?''
\hi{2}``Which piles do you want to hand over? Start with 1.''
\hi{3}``Compare with Problem 1.''}
\prow{\textbf{P11} \textit{(further)} Last counter loses; colour 0--20 for 1 or 2, and 1, 2 or 3.}
{1 or 2: 1, 4, 7, 10, 13, 16, 19. 1, 2 or 3: 1, 5, 9, 13, 17. Square 0 is not
coloured.}
{\hi{1}``Is 0 a good square to leave now?''
\hi{2}``Compare with Problem 2. What moved?''}
\prow{\textbf{P12} \textit{(further)} Two piles, 1 or 2 from one pile: 7 and 7. Explain.}
{\textbf{Second}: copy the opponent's move in the other pile. After each of
your moves the piles are equal, so you always have a reply and take the last
counter.}
{\hi{1}``Play 2 and 2, then 3 and 3.''
\hi{2}``What if you do what your opponent just did, in the other pile?''
\hi{3}``After your move the piles are equal. Can your opponent take the
last counter?''}
\prow{\textbf{P13} \textit{(further)} Colour the 0--9 chart for the P12 game; then 20 and 11.}
{Colour the 34 squares where the two piles have the same remainder on dividing
by 3 (the diagonal and the lines 3, 6 and 9 off it). 20 and 11 both leave 2:
\textbf{second}. Why: as with one pile, only the remainders matter. When they
match, any move changes one; take 1 or 2 from the pile with the bigger
remainder to make them match again. Note 1 and 4 is coloured though unequal.}
{\hi{1}``Fill row 0 and column 0 first: those are one-pile games.''
\hi{2}``Is 1 and 4 a good pair to leave? Play it.''
\hi{3}``Describe the coloured squares in words; use the words for 20 and
11.''}
\prow{\textbf{P14} \textit{(further)} Kim: for any list of moves the colouring repeats. Track for 2,
5 or 6.}
{2, 5 or 6: 0, 1, 4, 8, 11, 12, 15, 19, 22, 23, 26, 30, repeating every 11
(11--16 look like 0--5). \textbf{Kim is right} for a finite list: a square
depends only on the 6 squares below it (6 is the largest move); 6 squares can
be coloured in only 64 ways, so two of the first 65 blocks of 6 match, and
from there everything repeats.}
{\hi{1}``Colour the track. Where does it start to repeat?''
\hi{2}``To decide square 20, which squares do you look at?''
\hi{3}``How many ways can 6 squares in a row be coloured?''}
\end{probtable}

\section{The mathematics behind the week}

\textbf{Outcomes.} Call a pile \emph{P} if the player about to move loses with
best play (the squares the pages ask children to colour) and \emph{N}
otherwise. Pile 0 is P. A pile is N if some move reaches a P pile, and P if
every move reaches an N pile (or no move is possible). Induction on the pile
size labels every pile and gives the strategy: from N, move to P; from P every
move goes to N, so the opponent never reaches P and never reaches 0. A
strategy has to answer every reply; Lena and Omar test exactly that.

\textbf{Moves 1 to $k$.} P = multiples of $k+1$. From a multiple, a move of 1
to $k$ leaves a non-multiple; from $q(k+1)+r$ with $1 \le r \le k$, take $r$.
In play, answer a take of $s$ with $k+1-s$.

\textbf{The other rules, by remainders.} With 1, 3 or 4, P = remainder 0 or 2
on dividing by 7: from remainder 0 or 2 the three moves give remainders 6, 4,
3 or 1, 6, 5; from remainders 1, 3, 4, 5, 6, take 1, 1, 4, 3, 4. Likewise 1 or
4 gives 0, 2 mod 5; 2 or 3 gives 0, 1 mod 5 (1 is P because no move is
possible); and 1, 3 or 5 gives the even piles, since every move changes
parity. Greedy is often wrong: with 1, 3 or 4, from 8 the move to 4 loses.

\textbf{Periodicity.} If the largest move is $s$, the label of $n \ge s$
depends only on the labels of $n-1, \dots, n-s$. There are at most $2^s$ such
windows, so two of the windows starting at $0, 1, \dots, 2^s$ coincide
(pigeonhole), and from the first of them on the labels repeat. For 1, 3 or 4
the window 7--10 equals 0--3, so the period is 7 from the start; for 2, 5 or 6
the period is 11, with P $\equiv$ 0, 1, 4, 8 (mod 11).

\textbf{Designing a rule.} The P piles are exactly the multiples of $m$ if and
only if the move set contains $1, \dots, m-1$ and no multiple of $m$. Each of
$1, \dots, m-1$ must be N, and 0 is its only P option; a move that is a
multiple of $m$ would join two P piles. Sufficiency is the argument for moves
1 to $k$.

\textbf{Last counter loses.} With moves 1 to $k$, P = remainder 1 on dividing
by $k+1$: the player facing 1 must take it, and the same complement strategy
aims at 1 instead of 0. (The shift by one is special to moves 1 to $k$.)

\textbf{Two piles.} Copying (Tweedledum and Tweedledee) shows equal piles are
P in any game of this kind. In general give each pile its Grundy value $g(n)$,
the smallest number not among the values of the piles you can move to; two
piles are P exactly when $g(a) = g(b)$, and several piles when the XOR of the
values is 0 (Sprague--Grundy). With 1 or 2, $g(n) = n \bmod 3$, which gives the
34 squares of the 4--5 chart and makes 20 and 11 P. With 1, 3 or 4, $g$ repeats
0, 1, 0, 1, 2, 3, 2, so 5 and 6 (values 3 and 2) is N and 4 and 6 (both 2) is P.
In K--1 P8, 1 and 4 is P though unequal.

\textbf{Where it goes next.} Nim (any number from one of several piles; P
exactly when the XOR of the pile sizes is 0) and Sprague--Grundy for sums of
games; octal games such as Kayles, where a move may split a pile, and Guy's
open conjecture that every finite octal game is eventually periodic; misère
play; Wythoff's game. Sources: Berlekamp, Conway and Guy, \emph{Winning Ways};
Siegel, \emph{Combinatorial Game Theory}.

\newpage
\section{What to write down after the session}

Fill in one block per table in the five minutes after the session, and keep
the children's pages that show a claim. For each claim, circle how far it got:
\textbf{tried} (played a few games), \textbf{guessed} (said it without
testing), \textbf{checked in some cases} (tested on several piles or
replies), or
\textbf{explained} (gave a reason that covers every pile or every reply).

Claims to listen for. K--1: ``go 2nd with 3 and 6''; ``land on 9, 6, 3'';
``the secret changes with 1, 2 or 3''; ``copy your partner''. 2--3: ``leave 2
or 7''; ``the coloured squares go 2 then 5''; ``take as many as you can''
(false). 4--5: ``it repeats every 7''; ``when the last counter loses, every
coloured square moves up 1''; ``every rule repeats''.

\newcommand{\st}{\footnotesize tried \enspace guessed \enspace checked in some cases \enspace explained}
\newcommand{\recordblock}[2]{%
\textbf{#1}\par
\vspace{2pt}
#2\par
\vspace{4pt}
\noindent\begin{tabularx}{\linewidth}{@{}Xl@{}}
\toprule
\hd What a child tried, said or claimed (who) & How far it got (circle) \\
\midrule
\rule{0pt}{17pt} & \st \\ \arrayrulecolor{rule}\midrule
\rule{0pt}{17pt} & \st \\ \midrule
\rule{0pt}{17pt} & \st \\ \arrayrulecolor{black}\bottomrule
\end{tabularx}\par
\vspace{6pt}
A drawing or claim worth keeping:\par
\vspace{18pt}\hrule\vspace{12pt}}

\recordblock{K--1 table}{Last problem reached by each pair: \quad Pair 1: \underline{\hspace{0.6in}} \quad Pair 2: \underline{\hspace{0.6in}} \quad (of 9)}

\recordblock{2--3 table}{Last problem reached by each pair: \quad Pair 1: \underline{\hspace{0.6in}} \quad Pair 2: \underline{\hspace{0.6in}} \quad (of 9)}

\recordblock{4--5 table}{Last problem reached by each child: \quad \underline{\hspace{0.6in}} \quad \underline{\hspace{0.6in}} \quad \underline{\hspace{0.6in}} \quad (of 14)}

\end{document}
